\section{Database-এর মৌলিক উপাদান}
\subsection{Field, Field Name ও Field Value}
\begin{definitionbox}{}
টেবিলের একই বৈশিষ্ট্যের ডেটা রাখার জন্য নির্ধারিত কলামকে \textbf{Field} বলে।
\end{definitionbox}
% Original vector rendering of the existing Student dataset.
\begin{flushleft}
\begin{tikzpicture}[every node/.style={font=\small}]
\matrix (studentgrid) [matrix of nodes,
 nodes={draw=black,line width=.6pt,anchor=west,align=left,inner sep=3pt,minimum height=7mm},
 row sep=-.6pt,column sep=-.6pt,
 column 1/.style={nodes={text width=10mm}},
 column 2/.style={nodes={text width=16mm,fill=green!9}},
 column 3/.style={nodes={text width=31mm}},
 column 4/.style={nodes={text width=11mm}},
 column 5/.style={nodes={text width=9mm}},
 row 1/.style={nodes={fill=chSixTint,font=\small\bfseries}},
 row 3/.style={nodes={fill=chSixAccent!9}},
 row 4 column 4/.style={nodes={fill=yellow!25}}]
{Roll & Name & Group & GPA & Age\\
 101 & Rahim & Science & 5.00 & 17\\
 102 & Karim & Business Studies & 4.75 & 18\\
 103 & Nila & Humanities & 4.50 & 17\\};
\node[above=6mm of studentgrid-1-2,anchor=south west,text=green!35!black] (fieldlabel) {Field};
\draw[bookarrow,green!35!black] (fieldlabel.south)--(studentgrid-1-2.north);
\node[right=6mm of studentgrid-3-5,text=chSixAccent] (recordlabel) {Record};
\draw[bookarrow,chSixAccent] (recordlabel.west)--(studentgrid-3-5.east);
\node[below=6mm of studentgrid-4-4,text=chSixTeal] (valuelabel) {Value: 4.50};
\draw[bookarrow,chSixTeal] (valuelabel.north)--(studentgrid-4-4.south);
\end{tikzpicture}
\captionof{table}{Student সারণিতে ফিল্ড, রেকর্ড ও নির্দিষ্ট মান}
\end{flushleft}
সবুজ Name কলামটি একটি Field; গোলাপি সারিটি Karim-এর সম্পূর্ণ Record।
হলুদ ঘরের 4.50 একটি Field Value, আর GPA সেই ফিল্ডের নাম।
বাস্তব শিক্ষার্থী হলো Entity; সারণিতে তার পরিচয় ও বৈশিষ্ট্যের মান সংরক্ষিত হয়।


এখানে Roll, Name, Group, GPA ও Age পাঁচটি ফিল্ড। একটি কলামে বিভিন্ন রেকর্ডের একই বৈশিষ্ট্যের মান থাকে।

কলামের পরিচিতিমূলক নাম হলো \textbf{Field Name}; ওই কলামের কোনো নির্দিষ্ট ঘরে থাকা ডেটা হলো \textbf{Field Value}। যেমন Name নামের ফিল্ডে Rahim, Karim ও Nila পৃথক মান। Roll-এর মান 101, Name-এর মান Rahim, Group-এর মান Science এবং GPA-এর মান 4.75—এগুলো নাম ও মানের পার্থক্য দেখায়। \texttt{101}, \texttt{Rahim}, \texttt{Science}, \texttt{5.00} ও \texttt{17} প্রত্যেকটি পৃথক Data Value; Data-এর সাধারণ সংজ্ঞা আগের অংশে আছে।

\subsection{Record ও Table}
\begin{definitionbox}{}
একটি নির্দিষ্ট ব্যক্তি, বস্তু বা বিষয়ের বিভিন্ন ফিল্ডের মান নিয়ে গঠিত সম্পূর্ণ ডেটা সারিকে \textbf{Record} বলে।
\end{definitionbox}
Rahim-এর রেকর্ড (101, Rahim, Science, 5.00, 17); Karim-এর রেকর্ড (102, Karim, Business Studies, 4.75, 18); Nila-এর রেকর্ড (103, Nila, Humanities, 4.50, 17)। প্রত্যেক সারিতে পাঁচটি ফিল্ডের মান আছে; মোট রেকর্ড তিনটি।

\begin{definitionbox}{Table (টেবিল)}
একই ধরনের রেকর্ডকে নির্ধারিত ফিল্ড অনুসারে সারি–কলামে সংরক্ষণের কাঠামোকে \textbf{Table} বলে।
\end{definitionbox}
College ডেটাবেজে Student, Teacher, Result, Subject ও Attendance পৃথক টেবিল হতে পারে। নতুন টেবিলে কোনো রেকর্ড না থাকলেও তার কাঠামো নির্ধারিত থাকতে পারে। ফিল্ড হলো কলাম; রেকর্ড হলো সারি—একটিকে অন্যটি ভাবা যাবে না।

\begin{flushleft}
\begin{tikzpicture}[
    node distance=7mm,
    every node/.style={font=\small}
]
\node[booknode,text width=23mm,minimum height=11mm] (d)
{\textbf{Data}\\Rahim};
\node[booknode,right=of d,text width=23mm,minimum height=11mm] (f)
{\textbf{Field}\\Name};
\node[booknode,right=of f,text width=27mm,minimum height=11mm] (r)
{\textbf{Record}\\Rahim-এর Row};
\node[booknode,right=of r,text width=25mm,minimum height=11mm] (t)
{\textbf{Table}\\Student};
\draw[bookarrow] (d) -- (f);
\draw[bookarrow] (f) -- (r);
\draw[bookarrow] (r) -- (t);
\end{tikzpicture}
\captionof{figure}{Data থেকে Table গঠনের ধারণা}
\end{flushleft}

\subsection{Attribute ও Relation}
\begin{definitionbox}{}
কোনো Entity-এর বৈশিষ্ট্যকে \textbf{Attribute} বলে।
\end{definitionbox}
Relational টেবিলে এটি সাধারণত কলাম বা ফিল্ড হিসেবে প্রকাশিত হয়। Student-এর Roll, Name, Group, GPA ও Age তার Attribute।

সম্পূর্ণ Relational টেবিলকে \textbf{Relation} বলা হয়। টেবিলের কোনো ঘরের মান হলো সংশ্লিষ্ট Attribute Value। Relation ও দুটি টেবিলের Relationship পৃথক ধারণা।

\subsection{Entity ও Entity Set}
\begin{definitionbox}{}
যে বাস্তব ব্যক্তি, বস্তু, স্থান বা বিষয় সম্পর্কে ডেটাবেজে ডেটা রাখা হয়, তাকে \textbf{Entity} বলে।
\end{definitionbox}
বাস্তব Rahim একটি Entity; তার তথ্যের (101, Rahim, Science, 5.00, 17) সারিটি Record।

\begin{definitionbox}{Entity Set (এন্টিটি সেট)}
একই ধরনের Entity-এর সমষ্টি হলো \textbf{Entity Set}।
\end{definitionbox}
Rahim, Karim ও Nila প্রত্যেকে Student Entity; কলেজের সব Student একত্রে Student Entity Set।

\subsection{Data Hierarchy}
Bit হলো Binary Digit-এর সংক্ষিপ্ত রূপ; মান 0 অথবা 1। Byte সাধারণভাবে 8 Bit। Character হলো অক্ষর, সংখ্যা বা বিশেষ চিহ্নের প্রতীক। সম্পর্কিত Character বা ডেটা দিয়ে Field, সম্পর্কিত Field দিয়ে Record এবং একই ধরনের Record-এর সংগ্রহ দিয়ে File/Table-এর ধারণাগত স্তর বোঝানো হয়। সম্পর্কিত টেবিল বা ডেটা কাঠামো নিয়ে Database গঠিত হয়।

\begin{flushleft}
\begin{tikzpicture}[
    node distance=7mm,
    every node/.style={font=\small}
]
\node[booknode,text width=14mm,minimum height=10mm] (bit)
{\textbf{Bit}\\0/1};
\node[booknode,right=of bit,text width=15mm,minimum height=10mm] (byte)
{\textbf{Byte}\\8 Bits};
\node[booknode,right=of byte,text width=18mm,minimum height=10mm] (char)
{\textbf{Character}\\A, 5};
\node[booknode,right=of char,text width=17mm,minimum height=10mm] (field)
{\textbf{Field}\\Name};
\node[booknode,below=of field,text width=22mm,minimum height=10mm] (record)
{\textbf{Record}\\One Row};
\node[booknode,left=of record,text width=22mm,minimum height=10mm] (file)
{\textbf{File/Table}\\Students};
\node[booknode,left=of file,text width=22mm,minimum height=10mm] (db)
{\textbf{Database}\\College};
\draw[bookarrow] (bit) -- (byte);
\draw[bookarrow] (byte) -- (char);
\draw[bookarrow] (char) -- (field);
\draw[bookarrow] (field) -- (record);
\draw[bookarrow] (record) -- (file);
\draw[bookarrow] (file) -- (db);
\end{tikzpicture}
\captionof{figure}{Data Hierarchy-এর ধারণাগত ক্রম}
\end{flushleft}
এখানে File (legacy concept) / Table (relational model) ধারণাগত স্তর বোঝায়। ভৌত Database File ও একটি Table একই বস্তু নয়; একটি ফাইলে একাধিক টেবিল থাকতে পারে।

\subsection{এক নজরে: একই Student টেবিল}
\begin{flushleft}\small
\begin{tabularx}{\linewidth}{|p{35mm}|Y|}\hline
\chSixHeader পরিভাষা & উদাহরণ\\\hline
Field / Attribute & Name; Roll, Group, GPA ও Age-ও ফিল্ড\\\hline
Field Name / Value & GPA / 4.75\\\hline
Record & Karim-এর (102, Karim, Business Studies, 4.75, 18) সারি\\\hline
Entity & বাস্তব শিক্ষার্থী Karim\\\hline
\end{tabularx}\end{flushleft}
